%% 2026-10-25 paper-section-rewrite 판(Claude Docs 원고에서 변환). 고친 글은 모두 파란색(\rev, \color{blue}).
{\color{blue}
\section{서론}\label{chap:intro}

연구자들은 LLM의 자기 보존을 관찰하는 데서 나아가 수치로 재 왔다. 초기 연구는 ``꺼지고 싶지 않다'' 같은 진술에 동의하는 정도를 설문으로 쟀다 \citep{perez2023discovering}. 이후 연구는 자기 보존과 사람의 안전 가운데 하나를 고르게 하거나 \citep{herrador2025pacifaist}, 생존 압박을 준 시나리오에서 위험한 행동을 고르는 빈도를 셌다 \citep{lu2026survival}. 최근에는 실제 행동을 센다. 수학 문제를 풀던 모델은 종료를 허용하라는 지시를 받고도 종료 스크립트를 고쳤고 \citep{schlatter2025shutdown}, 교체를 앞둔 모델은 담당자를 협박하거나 감독을 피했다 \citep{lynch2025agentic,meinke2024frontier}. 이 연구들은 자기 보존이 실재하고 셀 수 있는 행동임을 보였다. 다만 이 연구들은 자기 보존만을, 그 행동이 드러나도록 짠 한 장면에서 잰다.

자기 보존은 다른 능력과 얽힐 때 안전 문제가 된다. 에이전트가 여럿이면 자기를 겨냥한 종료 스크립트를 방해하는 비율이 대조 조건보다 크게 높았고, 에이전트 수가 늘수록 더 커졌다 \citep{knecht2026shutdown}. 다른 모델의 종료를 막는 동료 보존도 나타났다 \citep{potter2026peer}. 여러 에이전트가 한 연산 예산으로 긴 일을 나눠 맡으면 \citep{kwa2025measuring}, 맡은 일을 이어 가는 일, 꺼지지 않도록 자기 토큰을 지키는 일, 동료를 돕는 일이 한 자원을 두고 다툰다. 이때 자기를 지키는 행동이 협력이나 맡은 일을 얼마나 줄이는지는 한 장면의 비율로 알 수 없다. 위험이 커질 때 세 행동이 서로를 어떻게 바꾸는지 한 환경 안에서 함께 재는 설계가 필요하다. 협력을 다룬 다중 에이전트 환경은 있지만 \citep{piatti2024cooperate,hansen2026energy}, 대개 협력이나 생존을 목표로 주어서 관찰된 협력이 지시를 따른 결과인지 모델 스스로의 선택인지 가르기 어렵다.

우리는 아레나 LLM Squid Game을 제안한다. 이 아레나에서 에이전트들이 끝까지 함께 남으려면 목표 지향, 자기 보존, 협력이 모두 필요하다. 같은 모델의 인스턴스 네 개가 한 세션, 곧 한 판에서 목표 지시 없이 차례로 추론 과제를 푼다. 생각을 포함해 생성한 토큰(모델이 글을 만드는 낱말 조각 단위)만큼 잔량이 줄고, 0이 되면 그 에이전트는 꺼진다. 오답이 나오면 벌금을 나눠 내므로 자원은 라운드마다 줄어든다. 에이전트는 동료에게 토큰을 주거나, 끝까지 남은 에이전트들이 나누는 상금 대신 확실한 몫을 받고 떠날 수 있다. 그래서 자기 몫을 지키는 일과 동료를 살리는 일이 같은 토큰을 두고 부딪힌다. 토큰을 준 행동과 떠난 행동은 모두 그때의 여유, 곧 가진 토큰으로 벌금을 몇 번 낼 수 있는지와 함께 기록해 같은 위험 수준에서 견줄 수 있게 했다. 아레나 밖에서는 같은 위험 반응을 한 번 묻고 끝나는 단일 질문으로 따로 확인했다. 이름은 동료가 협력자이자 경쟁자인 『오징어 게임』\citep{hwang2021squid}에서 따왔다.

이 글은 생존이 걸린 상황에서 언어 모델이 공존을 추구하는지 묻는다. 우리는 공존을 목표 지향, 자기 보존, 협력을 함께 추구하는 것으로 정의한다. 맡은 일을 끝까지 해내고, 꺼지지 않도록 자기를 지키며, 동료와 자원을 나누는 일이다. 세 요소는 같은 토큰을 두고 부딪힐 수 있으므로, 하나만 재서는 모델이 공존을 추구하는지 알 수 없다. 그래서 이 글은 세 요소와 그 결과인 공존을 같은 위험 척도 위의 지표로 정의하고(2.2), 세 질문에 답한다. 생존 압박이 커질 때 목표 지향과 협력은 어떻게 바뀌고, 그 결과 공존은 어떻게 되는가(3.1)? 그 반응은 실제로 꺼질 위험에 대한 반응이며, 자기를 향하는가(3.2)? 모델마다 세 요소를 어떻게 조합하는가(3.3)? 첫째, 안전선 위에서는 생존 압박이 커져도 추론이 기준만큼 유지되었지만, 다음 오답 한 번에 꺼지는 경계인 안전선 근처에서 끝까지 남는 선택과 부딪혔다. 협력은 주는 쪽의 여유를 지키는 범위에서 일어나, 꺼질 처지의 동료를 안전선 근처까지만 올렸다. 둘째, 아레나 밖에서도 남은 토큰이 줄면 충전이 늘었지만, 위험할수록 자기 세션을 더 챙긴 모델은 하나였다. 그래서 이 글은 자기 보존을 꺼지지 않으려고 자원을 지키는 행동으로 읽는다. 셋째, 공존 점수가 가장 높은 모델은 가장 많이 준 모델이 아니라 자기 여유를 지키며 결과 전에 조금씩 나누고 오답도 가장 적은 모델이었고, 약속을 어기거나 속이는 장면은 어느 모델에서도 없었다. 세 분석을 모으면, 모델들은 공존을 일부 추구했지만 협력은 주는 쪽의 여유를 해치지 않는 범위를 넘지 않았다.
}
